%\myChapterStar{Titre}{Titre court}{Ajouter à la table des matières? (false|true|chapter|section|subsection|subsubsection -chapter par défaut-)}
\myChapterStar{Introduction générale}{}{true}
%\myMinitoc{Profondeur de la minitoc (section|subsection|subsubsection)}{Titre de la minitoc}
\myMiniToc{}{Sommaire}

% \newpage
%\mySectionStar{Titre}{Titre court}{Ajouter à la table des matières? (false|true|chapter|section|subsection|subsubsection -section par défaut-)}
\mySectionStar{Contexte du travail}{}{true}
Une part importante des services proposés aujourd'hui sur le web repose sur des activités organisées sous forme de processus métier. Lorsqu'un usager effectue une inscription en ligne, dépose une demande administrative, suit une réclamation, réalise une commande ou valide une opération interne dans une organisation, il participe à une suite d'actions coordonnées, portées à la fois par des humains, des règles métier et des systèmes logiciels. Dans ce contexte, l'application web ne se limite plus à présenter des pages ou à stocker des informations : elle devient le support d'exécution d'un travail métier, avec des rôles, des formulaires, des données, des validations et des enchaînements d'activités.

%Cette réalité est encore plus visible lorsque l'on observe une organisation dans son ensemble. Une application métier ne prend généralement pas en charge une seule procédure isolée ; elle doit plutôt faire coexister plusieurs processus liés par des acteurs, des données et des décisions communes. Une commande peut modifier l'état du stock, un stock insuffisant peut déclencher un réapprovisionnement, une réclamation peut se rattacher à une vente déjà traitée. La difficulté n'est donc pas uniquement de transformer un diagramme en pages web, mais de produire une application capable de garder une cohérence entre plusieurs workflows appartenant au même contexte métier.

Le Business Process Management (BPM) fournit un cadre pour identifier, modéliser, analyser et améliorer ces processus. Les travaux de référence dans ce domaine montrent que la maîtrise d'un processus passe autant par sa représentation explicite que par sa capacité à être exécuté et adapté dans le système d'information \cite{dumas2018fundamentals,weske2012bpm}. Parmi les notations utilisées, BPMN occupe une place importante, car elle permet de décrire les acteurs, les événements, les tâches et les flux de contrôle d'un processus selon une notation largement reconnue \cite{omg2013bpmn}. Cependant, un modèle BPMN, même correct du point de vue du processus, ne suffit généralement pas à spécifier directement un socle applicatif web métier. Il indique ce qui doit s'enchaîner, mais ne décrit pas toujours avec précision les écrans attendus, les données manipulées, les champs visibles ou obligatoires, les règles de création des formulaires, ni la manière dont les utilisateurs interagissent concrètement avec le processus.

Parallèlement, l'ingénierie dirigée par les modèles propose de placer les modèles au centre du développement logiciel, afin de transformer des descriptions abstraites en artefacts plus proches de l'exécution \cite{omg2014mda,jezequel2012}. Cette orientation est particulièrement intéressante pour les applications web orientées processus, car elle permet d'envisager une continuité entre la compréhension métier, la spécification du processus et la production de l'application. Toutefois, les approches existantes restent souvent limitées à la génération de fragments d'interface, de prototypes ou de structures partielles. Plusieurs travaux ont exploré le passage de BPMN vers des modèles d'interaction ou vers des applications guidées par les processus \cite{diaz2018,sajji2022,sajji2023,drave2022}, mais le problème demeure lorsqu'il s'agit de produire un socle applicatif web métier structuré et cohérent avec les exigences exprimées dans le modèle.

C'est dans cette perspective que s'inscrit ce mémoire, intitulé « Ingénierie des applications web pilotées par des modèles BPMN ». Le travail présenté porte sur le framework ProOWeb, qui vise à générer un socle applicatif web métier à partir de modèles BPMN enrichis. L'idée centrale n'est pas seulement de représenter un processus ni de produire une maquette, mais de faire du modèle enrichi une source exploitable pour construire un socle capable de prendre en charge les acteurs, les tâches, les données, les formulaires, les transitions et le suivi des instances. Le framework ProOWeb cherche ainsi à rapprocher la modélisation métier et la réalisation logicielle, tout en conservant une organisation suffisamment claire pour permettre l'évolution du socle généré.

\mySectionStar{Problématique et objectifs}{}{true}
La difficulté principale vient du décalage entre la richesse attendue d'une application web métier et les informations habituellement contenues dans un modèle BPMN standard. Un processus peut préciser qu'une tâche doit être exécutée par un acteur donné, mais l'application correspondante doit aussi savoir quelles données afficher, quelles informations demander, quelles validations appliquer, quelles pages produire, comment relier les actions de l'utilisateur à l'avancement du processus et comment préserver les parties personnalisées lorsque le modèle évolue. Sans mécanisme structuré, le passage du modèle au code devient manuel, répétitif et fragile ; il expose le projet à des incohérences entre le processus conçu et l'application réellement livrée.

La problématique de ce travail peut donc être formulée de la manière suivante : comment concevoir un framework permettant, à partir de modèles BPMN enrichis, de générer un socle applicatif web métier fonctionnel, traçable et extensible, tout en conservant la cohérence entre les processus, les données, les rôles, les interfaces et l'évolution future de l'application ?

L'objectif principal de ce mémoire est de concevoir et de mettre à l'épreuve ProOWeb, un framework de génération d'applications web métier pilotées par des modèles BPMN enrichis. La réponse proposée à la problématique repose sur une chaîne explicite : le BPMN décrit le déroulement du travail, \textit{spec-v1} apporte les informations applicatives manquantes, une validation croisée contrôle leur cohérence, puis des contrats intermédiaires stabilisent les transformations modèle-à-modèle et modèle-à-texte. Ces transformations produisent un socle comprenant le backend, le frontend, les structures de données, les règles de sécurité et les ressources du runtime. Un manifeste relie enfin les sources aux fichiers générés et prépare leur régénération contrôlée.

De manière spécifique, il s'agit de :
\begin{itemize}
    \item étudier les notions de processus métier, de BPMN, d'applications web orientées processus et d'ingénierie dirigée par les modèles, afin de poser les bases conceptuelles du travail ;
    \item analyser les travaux existants sur la génération d'applications web à partir de modèles de processus, en mettant en évidence leurs apports et leurs limites ;
    \item définir et formaliser un modèle d'entrée enrichi associant BPMN et \textit{spec-v1}, ainsi que les contrôles nécessaires à leur mise en cohérence ;
    \item proposer l'architecture de ProOWeb, ses contrats intermédiaires, les algorithmes de transformation et ses mécanismes de traçabilité et de régénération ;
    \item mettre en œuvre cette chaîne et l'évaluer, dans les limites d'un cas d'étude, sur trois processus d'une même organisation.
\end{itemize}

La contribution ne se réduit donc pas à un outil de production de fichiers. Elle porte sur une démarche d'ingénierie qui précise les informations attendues, contrôle les modèles d'entrée, organise leur transformation et conserve la trace des décisions projetées dans le code. ProOWeb vise ainsi un socle applicatif fonctionnel, traçable et extensible, que le développeur doit encore vérifier et adapter aux contraintes propres à son organisation.

\mySectionStar{Organisation du manuscrit}{}{true}
La suite de ce mémoire est organisée en quatre chapitres, suivis d'une conclusion générale.

\textbf{Chapitre I — Généralités sur le Business Process Model and Notation et les applications web métier.} Ce chapitre présente les concepts de base nécessaires à la compréhension du sujet. Il revient sur les processus métier, le BPM, la notation BPMN, l'ingénierie dirigée par les modèles et les applications web métier. Il permet ainsi de fixer le vocabulaire et le cadre conceptuel utilisés dans la suite du mémoire.

\textbf{Chapitre II — Revue de la littérature sur la génération d'applications web pilotées par des modèles.} Ce chapitre étudie les travaux existants liés à l'automatisation des processus, aux systèmes BPM, à la génération de code, aux approches low-code/no-code et aux transformations de modèles. Il met en évidence les limites rencontrées lors du passage d'un modèle de processus à un socle applicatif web métier fonctionnel et évolutif.

\textbf{Chapitre III — Framework d'ingénierie proposé pour la génération d'applications web pilotées par BPMN.} Ce chapitre constitue le cœur de la contribution. Il présente ProOWeb, son principe général, son architecture, son modèle d'entrée enrichi, ses règles de transformation et les mécanismes prévus pour produire un socle applicatif web métier, que le développeur peut ensuite vérifier, adapter et compléter.

\textbf{Chapitre IV — Mise en œuvre et expérimentation.} Ce chapitre est consacré à la réalisation concrète du framework et à son évaluation. Il présente le socle généré pour une organisation de commerce électronique, les choix de mise en œuvre, les trois processus retenus, les interfaces obtenues, les résultats observés ainsi que les limites du travail réalisé.

\textbf{Conclusion générale —} elle rappelle la problématique traitée, synthétise les contributions du mémoire, discute les résultats obtenus et ouvre quelques perspectives pour l'amélioration de ProOWeb.

\myCleanStarChapterEnd
